\section{\texttt{TLOG}}
\noindent\textbf{Descriptor profile.} Vec entry 24, tile catalog opcode \texttt{0x1030}. Descriptor bundle: \texttt{B.OP + B.IOT}; WO/WM sites: 3; VEC beat-level sites: 3; ROM bytes: 24.\par\smallskip
\TileInstructionEncoding{figs/encoding/tile_instruction_entries/entry_24_pto_tlog_th0.pdf}{figs/encoding/tile_instruction_entries/entry_24_pto_tlog_th2.pdf}{0x1030}{24}{B.OP + B.IOT}
\paragraph{Bundled descriptor (PTO-ISA header).}
\begin{description}[leftmargin=0.18\linewidth,style=nextline]
\item[Bundle] \texttt{B.OP + B.IOT}. \texttt{B.OP.desc\_count} = 1; \texttt{B.OP.rom\_entry\_id} = 24.
\item[Assembly] 0: \texttt{B.OP} selects \texttt{TLOG}, carries element format/variant, declares \texttt{desc\_count}, and names the ROM entry. 1: \texttt{B.IOT} binds architectural tile operands. Across all \texttt{B.IOT} descriptors: tile inputs \texttt{src}, mask inputs none, and outputs \texttt{dst}.
\item[Tile operands] 1 tile input(s): \texttt{src}; 1 tile output(s): \texttt{dst}. Bound by 1 architectural \texttt{B.IOT} descriptor; each \texttt{B.IOT} supplies at most two source selectors plus one destination selector.
\item[Scalar GPR operands] No scalar GPR import/export descriptor is required.
\item[Metadata operands] No additional dimension or attribute descriptor is required beyond the tile operand bindings.
\end{description}
\par\smallskip
{\small
\paragraph{PTO ISA description.}
Elementwise natural logarithm of a tile.
\paragraph{PTO ISA operation.}
For each element \texttt{(i, j)} in the valid region:

\[
\mathrm{dst}_{i,j} = \log(\mathrm{src}_{i,j})
\]
}\par\smallskip
\paragraph{Microcode site table.}
{\scriptsize
\begin{longtable}{@{}R{0.055\linewidth}L{0.23\linewidth}L{0.31\linewidth}L{0.22\linewidth}@{}}
\toprule
Seq & Micro instruction & WO[63:32] fields & WM[31:0] fields \\
\midrule
0 & \texttt{u\_vec.vlds} & \texttt{opc=0x17, x=0, fl=mem, seq=0, kind=b} & \texttt{entry=24, cnt=2, res=vec, var=0} \\
1 & \texttt{u\_vec.vln} & \texttt{opc=0x18, x=0, fl=alu, seq=1, kind=u} & \texttt{entry=24, cnt=1, res=vec, var=0} \\
2 & \texttt{u\_vec.vsts} & \texttt{opc=0x2A, x=0, fl=mem, seq=2, kind=b} & \texttt{entry=24, cnt=2, res=vec, var=0} \\
\bottomrule
\end{longtable}
}
\paragraph{ROM body DSL.}
\begin{lstlisting}[style=ptocode]
def rom_tlog(src, dst):
    dtype = dst.element_type
    valid_rows, valid_cols = dst.valid_shape

    for row in range(0, valid_rows, 1):
        remained = valid_cols
        for col in range(0, valid_cols, pto.get_lanes(dtype)):
            mask, remained = pto.make_mask(dtype, remained)
            vinput = pto.vlds(src[row, col:])
            result = pto.vln(vinput, mask)
            pto.vsts(result, dst[row, col:], mask)
    return
\end{lstlisting}
\Needspace{0.34\textheight}
